\section{Business Opportunity}

The freelance market has grown faster than the infrastructure that makes it safe. Most platforms solved the discovery problem years ago: a client can publish a job in minutes and receive proposals within days. What remains unsolved is the settlement problem. A client still has to transfer money to someone whose identity is asserted rather than demonstrated, and a freelancer still has to start work on the promise of a future payment. The cost of that uncertainty is paid in every transaction, through higher prices for clients, delayed payment for freelancers, and a large volume of work that never leaves the platform's support queue because there is no shared record of what was agreed.

This creates three concrete openings that FREVIA is designed to occupy.

The first is verifiable identity. A marketplace that keeps evidence of an identity check, together with the result of an AI document and liveness pipeline and an administrator's review, can let a freelancer carry a badge that means something. This is especially relevant for cross-border engagements, where the two parties share no legal jurisdiction, no mutual contact, and no practical way to verify each other before work begins. FREVIA stores the verification outcome and its supporting documents, exposes only the resulting status to other users, and keeps the documents accessible to administrators and to the document's owner. The same mechanism can later support compliance obligations that a marketplace operator may be required to demonstrate.

The second is escrow at a granular level. Clients are far more willing to commit money when the commitment is bounded. FREVIA asks the client to fund one milestone at a time rather than the whole contract, and it publishes the resulting payment status -- pending, funded, released, refunded, or disputed -- on the milestone itself. A freelancer can therefore see that the money for the current milestone has already been collected before starting work, while the client retains the ability to review the deliverable and to request changes with a stated due date before releasing funds. This structure converts a single high-stakes payment decision into a series of low-stakes ones, which is precisely the change that makes a first engagement between strangers possible.

The third is a dispute procedure that both parties can operate. On most platforms a dispute disappears into customer support, and neither party can predict the outcome, the timeline, or the evidence that will be considered. FREVIA makes the procedure explicit and symmetric: the dispute is attached to a milestone, both parties owe a response within a stated deadline, both may upload evidence tagged as a claim or a response, an arbitration fee is charged with its own payment deadline, and the administrator's decision records a split of the milestone amount together with a written reason. The losing party can request a review, which produces a final decision. Because the whole sequence is stored, the outcome is defensible after the fact, and the threat of a well-documented dispute discourages the behaviour that makes disputes common in the first place.

Reliability is not only a matter of money. A large share of collaboration failure comes from communication that is scattered across private channels. FREVIA keeps the conversation inside the contract: one-to-one real-time messaging with attachments, read receipts and typing indicators, and shared files attached to the contract rather than to a chat thread. A persistent notification record ensures that a proposal decision, a milestone submission, a payment event, or a dispute deadline is not lost because nobody opened a particular page. A community forum, with an AI moderation step and mandatory administrative approval before publication, adds a public layer where users can ask questions and compare experiences.

The AI strategy is deliberately narrow, and that narrowness is itself a commercial position. FREVIA uses AI where the output is verifiable and reversible: extracting candidate skills from an uploaded CV so that a freelancer does not have to type them, verifying an identity document and a live face so that an administrator does not have to inspect every upload, and scoring forum content so that a moderator's queue is ordered by risk. It does not use AI to decide who should be hired. That decision is left to the client, with the platform supplying structured, comparable information -- skills with proficiency levels, portfolio items, verified identity status, completed contracts, and ratings -- instead of an opaque score. For a marketplace that is still building trust with its first users, an explainable process is worth more than an automated one.

The bilingual requirement reinforces the same positioning. FREVIA ships in English and Vietnamese, and the language is a user preference stored in a cookie rather than part of the URL, so switching language does not break a shared link to a job or a contract. For a platform launched from Vietnam and aimed at freelancers who work for international clients, this is not a localisation afterthought but a primary adoption path: the same users need to read job descriptions in English and to operate the contract, milestone, and dispute screens in Vietnamese.

Taken together, these choices define a market position that is narrower than Upwork's and deeper than a listing board's. FREVIA targets engagements where both parties need a defensible record: a verified identity on one side, funded escrow on the other, and a written procedure for the case where the two disagree. That segment is small enough to serve with a focused product and large enough to validate the model before any expansion into hourly contracts, multi-currency settlement, or additional verification methods.
